\chapter*{前言}
\addcontentsline{toc}{chapter}{前言}

作为最重要的自然科学之一，物理学涉及观测和理解两个层面，二者缺一不可。没有观测就没有人们试图去理解的事实或现象；而没有理解的观测充其量只能是某种记录而不能称其为科学。尽管经过了几百年的发展，物理学的使命从未改变，那就是在不同时空层次上去探究物质的组成和结构、它们之间的相互作用和演化动力学，以及研究对象的部分或整体在不同条件下所展现出的特性及其背后的物理规律。从亚原子层次上的新粒子及新物理，到原子与分子层次上气态、液态、固态、等离子态物质的性质与相变及其在外加条件下的新奇现象，再到宇观层次上星系动力学或宇宙的起源，各领域的物理学家们始终在践行着这一崇高使命：观测和理解我们所处的世界。

长时间以来，物理学的发展或多或少地受到了还原论的影响，即倾向于把复杂的系统分解成简单的各个部分，观察和理解各部分的性质和这些部分之间的相互作用，继而研究它们是如何融为整体并导致整个复杂系统所展现出的固有特性和动力学行为的。事实上，还原论的思想方法激励着现代科学不断发展，雄心勃勃的大统一理论就致力于通过研究微观粒子之间仅有的四种相互作用力（万有引力、电磁力、强相互作用力、弱相互作用力）之间的联系，寻找能够统一描述四种相互作用力的理论或模型。还原论是一种自上而下（top-down）的思维和研究范式。然而，很多复杂系统，例如凝聚相物质的宏观性质和相变动力学，以及生命的产生、生命个体或群落的行为等，并不一定能够通过分解为部分的叠加来得到正确理解。面对这些复杂系统，还原论的方法显得无能为力或者颇有局限性。科学家们发现，某些系统所展现的很多宏观性质和现象，是在系统的演化中逐步涌现的（emergent），因此一种自下而上（bottom-up）的涌生论研究范式应运而生。实际上，与涌生论相关的系统论，在控制论、信息论、运筹学、决策与博弈论、人工智能等学科中，也占有着十分重要的地位。在现代物理学的研究中，还原论与涌生论的范式相辅相成，在不同层次、不同学科的研究中均起着重要的推动作用。

为了理解选定的一个物理系统，无论是通过实验还是理论来研究，都需要操纵和控制某些作为输入的观测量，再去测量或者计算另外一些作为输出的观测量。然而，真实的世界过于复杂，为了理解所选定物理系统展现出的现象，就需要对真实系统进行简化，基于某些假设在数学上建立一定的模型。在能成功解释所观察到的现象的前提下，模型所需要的假设越少，这个模型就被认为是越优美的理论。物理学上的 Newton 运动定律、Maxwell 方程组、Schrödinger 方程等，无不展现出这种简洁而有力的美。然而，当这些理论用于真实系统时，能够精确求解的少之又少，我们就需要对这些理论做进一步的近似，这种近似可能是解析上的，也可能是数值上的。一方面，人们希望做适当的简化和近似，使得数学问题能够便于解析求解和理论分析；另一方面，随着超级计算机和人工智能技术的发展，人们有时倾向于从最一般的理论出发做从头计算，以便计算的结果更能反映真实的系统或者把系统描述得更加全面和准确。因此，在当今的自然科学中，数值计算已经与实验研究、理论分析形成了三足鼎立之势，相辅相成，共同推动着科技与文明的进步。此外，数值计算在经济、社会、人文与管理等学科中也占据着举足轻重的地位。值得注意的是，前面谈到的复杂系统，随着个体数目和变量的增加，如果按照第一性原理来处理，其所需的计算量是呈指数级增加的，会很快达到一个惊人而无法处理的地步。所幸的是，人们可以将复杂科学问题映射为高维参数空间的聚类、关联、反演、最优点搜寻等问题，并利用日益强大的人工智能技术来探索新的解决方案。通过定量发散与定性收敛相结合的研究方法，人工智能深度融合了先进算法、科学数据、智能计算系统等工具与第一性原理，将科学研究推进到了新的高度，且极大地提高了效率与产出，正引领着一次新的科学研究范式的转换。

无论是从基本物理原理出发的从头计算，还是利用计算机去模拟简化后的模型或者处理实验观测数据，数值计算都发挥着十分重要的作用。因此，对于学习物理学或者从事相关研究的工作者来说，深入了解和掌握数值计算的基本知识，并能正确运用于学习和工作之中，就显得尤为重要了。这里所说的“正确”二字，尤为值得再三强调。我们必须要熟悉数值计算的总体特点和处理各种数值问题的具体方法，成为一个聪明的软件使用者或程序开发者，时刻对数值模拟的过程和结果保持一个清醒的头脑。一个最重要的警告就是，虽然数值模拟的每一个子过程看似都是决定性和可预测的，然而数值模拟的结果却不一定是确定或者正确的。这种不确定性或者不正确性，首先可能源于过于众多的输入变量或对这些输入变量精确度的过分要求、模拟过程中的高度非线性系统、有限字长带来的舍入误差的不可控放大、算法设计本身等。如何甄别结果的正误和可靠性，就需要我们具有良好的素养和习惯。对于数值结果，要从不同角度进行探究，经过全方位的考察之后方可决定是否予以采用。例如：简化和抽象出的数学模型，是否能忠实描述真实的物理系统？模拟过程中的取样，是否具有足够的代表性？不同取样之间是否引入了人为的关联性？我们是否能用不同的数值模拟方法产生相同的结果？得到的结果与实验测量、基本理论预测或者前人的模拟结果，相比较如何？数值程序应用于极限情况下时，是否能产生理论模型在极限情况下应有的解析预测结果？你所得到的结果足够精确吗？实际上，无论是使用成熟的软件，还是自己编写程序，都应该养成谨慎考察和使用计算结果的良好习惯，不要轻率地相信自己的计算结果。

本书就是为了让物理学及相关专业的高年级本科生深入了解数值计算，用之解决常见物理问题，并养成良好的数值计算素养而编写的。实际上，国内外已经有很多优秀的计算物理学教材，其取材、体系和风格各不相同，有的偏重数值计算本身，有的偏重物理学中的个别分支学科的数值计算和物理讨论。但在编者看来，将数值计算方法与高等学校物理专业的主要课程（普通物理和四大力学等）有机融合起来，而且将二者的主要内容从易到难按照一定逻辑顺序进行介绍的教材，尚不多见。基于编者在北京大学物理学院主讲本科生主干基础课“计算物理学”十年来的教学经验，本书是在这方面的一个尝试。然而，为达此目的，在体系安排和内容取舍上是极具挑战性的。图1是本书涉及的数值分析基本内容及其逻辑关系图。在数值计算方法和涉及的计算物理学典型问题两个方面，我们都刻意避免面面俱到，而是尽量选择具有典型性和代表性的材料，力图形成一个有机的整体，使得同学们修完本课程后，能够对如何利用数值计算探讨和解决物理问题有一个整体的把握，并为进一步深造和未来的科研工作打下坚实的基础，培养必要的科学计算素养。

\begin{figure}[htbp]
\centering
\begin{tikzpicture}[
  x=1cm,y=1cm,>=Stealth,
  every node/.style={font=\scriptsize,align=center},
  box/.style={draw,rectangle,inner xsep=3pt,inner ysep=3pt,minimum height=6mm}
]
\node[box,minimum width=55mm] (top) at (7.0,8.1) {数值分析的主要内容};
\node[box,minimum width=24mm] (lin) at (3.3,6.9) {线性问题};
\node[box,minimum width=24mm] (nonlin) at (7.0,6.9) {非线性问题};
\node[box,minimum width=29mm] (rand) at (11.7,6.9) {随机问题\\与统计方法};

\node[box,minimum width=34mm] (linsolve) at (3.3,5.7) {线性方程组的求解};
\node[box,minimum width=25mm] (iter) at (7.0,5.7) {迭代方法};
\node[box,minimum width=24mm] (mc) at (11.7,5.1) {蒙特卡洛\\方法};

\node[box,minimum width=23mm] (fourier) at (0.6,4.35) {Fourier 分析};
\node[box,minimum width=23mm] (approx) at (3.3,4.35) {函数逼近};
\node[box,minimum width=25mm] (eig) at (6.0,4.35) {本征值问题};
\node[box,minimum width=30mm] (rootopt) at (9.1,4.35) {数值求根与\\非线性优化};

\node[box,minimum width=23mm] (diff) at (2.2,3.15) {数值微分};
\node[box,minimum width=23mm] (integ) at (4.9,3.15) {数值积分};
\node[box,minimum width=24mm] (markov) at (11.7,3.65) {\rusname{Марков}\\过程};

\node[box,minimum width=29mm] (ivp) at (2.2,1.75) {常微分方程的\\初值问题};
\node[box,minimum width=29mm] (bvp) at (6.4,1.75) {偏微分方程的\\边值问题};
\node[box,minimum width=29mm] (ml) at (11.7,1.65) {机器学习和\\人工智能};
\node[box,minimum width=29mm] (pde) at (4.9,0.35) {偏微分方程的\\数值求解};

% 主分支
\draw[->] (top.south west) -- ++(0,-0.25) -| (lin.north);
\draw[->] (top.south) -- (nonlin.north);
\draw[->] (top.south east) -- ++(0,-0.25) -| (rand.north);
\draw[->] (lin)--(linsolve);
\draw[->] (nonlin)--(iter);
\draw[->] (rand)--(mc);

% 线性/非线性数值方法之间的关系
\draw[->] (linsolve.east)--(iter.west);
\draw[->] (linsolve.south) -- ++(0,-0.35) -| (approx.north);
\draw[->] (linsolve.south) -- ++(0,-0.35) -| (eig.north);
\draw[->] (approx.west)--(fourier.east);
\draw[->] (iter.south) -- (rootopt.north);
\draw[->] (approx.south) -- (diff.north);
\draw[->] (approx.south) -- ++(0,-0.35) -| (integ.north);
\draw[->] (eig.south) -- ++(0,-0.35) -| (integ.north);
\draw[->] (eig.south) -- ++(0,-0.35) -| (bvp.north);
\draw[->] (diff)--(ivp);
\draw[->] (ivp.east) -- ++(0.45,0) |- (pde.west);
\draw[->] (integ.south) -- (pde.north);
\draw[->] (bvp.south) -- (pde.north east);

% 随机问题与统计方法分支
\draw[->] (mc)--(markov);
\draw[->] (markov)--(ml);
\draw[->] (rand.east) -- ++(0.45,0) |- (ml.east);
\end{tikzpicture}
\caption{本书力求涉及的数值分析基本内容及其逻辑关系图}
\end{figure}

本书适合用作双一流大学物理学专业及相关学科本科高年级课程的教材或教学参考书，对理工科研究生和教师，以及其他科技工作者也有一定的参考价值。

在教材出版过程中，国内许多专家对书稿提出了宝贵的建议和意见，其中有很多我们已尽力采纳修改。编者要特别感谢“101计划”审稿专家，北京理工大学姚裕贵教授、冯艳全教授，南开大学李宝会教授，北京航空航天大学周洪波教授，国防科技大学戴佳钰教授。

由于编者的水平有限，书中不当和疏漏之处在所难免，希望读者不吝赐教，将发现的问题或者您的建议和意见通过电子邮件或其他方式反馈给编者，以便再版时予以完善。我们的联系方式如下：

彭良友，邮件地址：liangyou.peng@pku.edu.cn，电话：010-62765027；梁昊，邮件地址：haoliang@pku.edu.cn；陈基，邮件地址：ji.chen@pku.edu.cn。
\clearpage
